\documentclass[11pt,a4paper]{article}

\usepackage{amsmath,amssymb,amsthm}
\usepackage{geometry}
\usepackage{booktabs}
\usepackage{natbib}
\usepackage{hyperref}

\geometry{margin=2.5cm}

\newtheorem{proposition}{Proposition}

\title{A stylised model of strategic intraday repositioning\\under
capacity, dual-priced imbalance settlement, and the ISP window}
\author{}
\date{}

\begin{document}

\maketitle

\begin{abstract}
A two-subperiod, two-market sketch illustrates how three Spanish regime
changes --- ISP60$\to$ISP15 (2024-12-01), MTU60$\to$MTU15 intraday
(2025-03-19) and DA60$\to$DA15 (2025-10-01) --- map to the sign and
magnitude of the dominant firm's net intraday position. Two primitives
do the work: a physical capacity and the Spanish dual-pricing rule. The
imbalance-settlement window (ISP60 vs ISP15) governs whether within-hour
imbalances can be netted before settlement, and therefore how strongly
the firm prices its structural day-ahead overcommitment into intraday
trade.
\end{abstract}

\section{Setup}

Two subperiods per hour, $t \in \{1,2\}$, with physical capacity
$K_t > 0$ and marginal cost $c$. Three stages.

\paragraph{DA.} The firm commits a block $q_{1t}^{\mathrm{phys}}$ at
price $p_1$. Under DA60, the block is uniform across subperiods:
$q_{11}^{\mathrm{phys}} = q_{12}^{\mathrm{phys}} = q_1/2$. Under DA15,
the two commitments are chosen independently.

\paragraph{Intraday.} The firm trades $q_{2t}$ at $p_{2t}$ against a
linear residual demand with slope $b_{2t}$:
$q_{2t}=b_{2t}(p_{1}-p_{2t})$.

\paragraph{Real time.} Let $q^s_t \equiv q_{1t}^{\mathrm{phys}} + q_{2t}$
be the firm's scheduled position. It produces $Q_t \in [0,K_t]$ and
generates an imbalance $\Delta_t = Q_t - q^s_t$, which is settled against
a random system state $Z\in\{H,L\}$, $\Pr(Z=L)=\pi$, by the Spanish
directional rule:
\[
p^{\mathrm{imb}} \;=\; \begin{cases}
p_1 & \text{deviation helpful to the system,} \\
p^{\mathrm{up}} & \text{harmful short } (\Delta<0,\,Z=L), \\
p^{\mathrm{dn}} & \text{harmful long } (\Delta>0,\,Z=H),
\end{cases}
\qquad p^{\mathrm{dn}} \leq p_1 \leq p^{\mathrm{up}}.
\]
Write the reserve spreads as
$\phi^+ \equiv p^{\mathrm{up}}-p_1$ and
$\phi^- \equiv p_1 - p^{\mathrm{dn}}$. Settlements aggregate over an ISP
window $\omega \in \{60,15\}$: under ISP60 the firm pays against the
hourly net $\Delta_H = \Delta_1 + \Delta_2$, so a slack slot cancels a
short slot; under ISP15 each $\Delta_t$ is settled independently, and no
within-hour offset is possible.

\section{Effective intraday marginal cost}

On the short branch the firm is paid at least $p_1$ per MWh produced, so
it produces up to capacity whenever its schedule $q^s_t$ exceeds $K_t$.
In that \emph{oversell} regime the imbalance is mechanical,
$\Delta_t = K_t - q^s_t < 0$, and the expected cost of being short pins
down the firm's effective intraday marginal cost.

Let $\kappa(\omega)\in[0,1]$ denote the share of the firm's hourly gross
short that is \emph{not} cancelled by within-hour offsets:
\begin{equation}
\kappa(\omega) \;=\;
\begin{cases}
\dfrac{|\Delta_1+\Delta_2|}{|\Delta_1|+|\Delta_2|} & \omega = 60, \\[8pt]
1 & \omega = 15.
\end{cases}
\label{eq:kappa}
\end{equation}
Under DA60, the block constraint forces $q_{11}^{\mathrm{phys}} =
q_{12}^{\mathrm{phys}}$, so whenever capacity varies within the hour
($K_1 \neq K_2$) the firm is short in the tight slot and slack in the
loose one. ISP60 lets it net; ISP15 does not.

\begin{proposition}[Effective intraday marginal cost]
\label{prop:mcstar}
On the oversell branch,
\[
\mathrm{MC}^\star(\omega) \;=\; p_1 + \pi\,\phi^+\,\kappa(\omega).
\]
On the undersell branch --- achieved under DA15 by the slot-by-slot match
$q_{1t}^{\mathrm{phys}} = K_t$ ---
\[
\mathrm{MC}^\star \;=\; p_1 - (1-\pi)\,\phi^-.
\]
\end{proposition}

\begin{proof}
On the oversell branch, $\Delta_t = K_t - q_t^s < 0$; producing $Q_t = K_t$ is
optimal since $p_1 > c$. Increasing the intraday sale $q_2$ by one unit raises
$q^s$ by one and deepens the short by one. A fraction $\kappa(\omega)$ of this
extra short is exposed at settlement; of that fraction the firm pays $\phi^+$
with probability $\pi$ (harmful-short state $Z=L$). The first-order condition
for the intraday sale therefore requires
\begin{equation*}
p_2 \;=\; p_1 + \pi\,\phi^+\,\kappa(\omega) \;\equiv\; \mathrm{MC}^\star(\omega).
\end{equation*}
On the undersell branch the firm sets $q_{1t}^{\mathrm{phys}} = K_t$ and
produces exactly its schedule, so $\Delta_t \geq 0$. A marginal increase in
$q^s$ displaces a unit that would otherwise be beneficial-long; it costs
$(1-\pi)\phi^-$ in expectation (harmful-long state $Z=H$), giving
$\mathrm{MC}^\star = p_1 - (1-\pi)\phi^-$.
\end{proof}

If $\phi^+=\phi^-=0$, then $\mathrm{MC}^\star \equiv p_1$ and imbalance
settlement drops out.

\section{Regime map}
\label{sec:regimes}

In the intraday stage the firm takes $p_1$ as given and maximises
profit against the linear residual demand of each subperiod.  For
subperiod $t$ the intraday clearing price is
$p_{2t} = p_1 - q_{2t}/b_{2t}$, so the first-order condition gives
$q_{2t} = \tfrac{b_{2t}}{2}(p_1 - \mathrm{MC}^\star)$.
Summing over $t\in\{1,2\}$ and writing $\beta\equiv b_{22}/b_{21}\in(0,1]$
for the within-hour residual-demand asymmetry and
$\Phi(\beta)\equiv(1+\beta)/2$:
\begin{equation}
q_{21}+q_{22} \;=\; b_{21}\,\Phi(\beta)\,\bigl(p_1 - \mathrm{MC}^\star\bigr).
\label{eq:xform}
\end{equation}
The four Spanish regimes map as follows.

\begin{center}
\begin{tabular}{llcc}
\toprule
Regime & Dates & $\mathrm{MC}^\star - p_1$ & $\mathrm{sign}(q_{21}+q_{22})$ \\
\midrule
R1 DA60/ID60, ISP60 & pre 2024-12-01      & $+\pi\phi^+\kappa_{60}$ & $-$ \\
R2 DA60/ID60, ISP15 & 2024-12-01 $\to$ 2025-03-19 & $+\pi\phi^+$ & $-$ \\
R3 DA60/ID15, ISP15 & 2025-03-19 $\to$ 2025-10-01 & $+\pi\phi^+$ & $-$ \\
R4 DA15/ID15, ISP15 & post 2025-10-01     & $-(1-\pi)\phi^-$        & $+$ \\
\bottomrule
\end{tabular}
\end{center}

\paragraph{The sign flip at R4.} Under DA15 the firm matches
$q_{1t}^{\mathrm{phys}}$ to $K_t$ slot-by-slot, closing the structural
oversell. It enters the undersell branch and
$\mathrm{MC}^\star < p_1$, so $q_{21}+q_{22}>0$. This flip is
closed-form: it needs only the mechanical removal of the DA block and
the dual-pricing rule, not a calibration of $\phi^\pm$ or $\pi$.

\paragraph{The ISP15 channel (R1$\to$R2): predicted but weak in the data.}
ISP15 sets $\kappa=1$ and strictly raises the imbalance cost of a given
DA overcommitment. In the model, firms internalise this at the DA stage:
the equilibrium $q_1$ under ISP15 is strictly lower than under ISP60,
shrinking the structural short and hence $|q_{21}+q_{22}|$
through~(\ref{eq:xform}). Empirically, however, this channel barely
fires at the 2024-12-01 break. Conditioning on low wind (where the
oversell branch should bind most tightly), the OMIE-classified dominant
firms' per-unit intraday repositioning $\Delta Q$ is essentially flat
across ISP60$\to$ISP15: $-268$ MWh/unit-day in the DA60/ID60 3-session
window (2024-06-14 $\to$ 2024-11-30) versus $-271$ MWh/unit-day in the
DA60/ID60 ISP15 window (2024-12-01 $\to$ 2025-03-19). The unconditional
and firm-level aggregates show the same pattern. One reading consistent
with the mechanism but with a second-order calibration: dominants' DA
bids are driven primarily by residual-demand curvature (market-power
rents) rather than by the expected imbalance penalty, so a sharper
penalty remains second-order as long as firms operate far from the
penalty-binding region. The penalty only binds once within-hour
repositioning is simultaneously restricted --- i.e.\ jointly with MTU15
(see R2$\to$R3 below).

\paragraph{The thinness channel (R2$\to$R3): the dominant empirical
compression.} MTU15 in the intraday market introduces asymmetry between
the two subperiods' residual-demand slopes; $\Phi(\beta)$ falls below
its $\beta=1$ value, further compressing $|q_{21}+q_{22}|$ without
flipping its sign. This is where the compression visible in the Spanish
data actually occurs: at the 2025-03-19 break, low-wind dominant
per-unit $\Delta Q$ falls sharply from $-271$ to $-78$ MWh/unit-day, a
71\% compression in a single step. The ordering holds unconditionally,
under firm-level aggregation using OMIE's \texttt{grupo\_empresarial}
labels (both an all-labeled-firms and a Big-4-only definition of
dominance), and under a balanced-panel restriction to firms present in
every regime, so composition effects and renewable buildout are ruled
out as primary drivers. The channel operates on two fronts at once:
MTU15 lowers $\Phi(\beta)$ through within-hour asymmetry in
residual-demand slopes, and it simultaneously activates the ISP15
penalty by removing the firm's ability to reposition at a coarser
granularity than the settlement grid. What the ISP60$\to$ISP15
transition failed to deliver in isolation (R1$\to$R2 essentially flat),
it delivers jointly with intraday granularity at R2$\to$R3.

\section{Summary}

Two primitives --- capacity and dual-priced settlement --- plus the ISP
window and the DA granularity generate the sign sequence $(-,-,-,+)$
across R1--R4, placing the ISP15 and DA15 reforms on distinct structural
channels: ISP15 compresses the magnitude through the DA-side response
to a higher effective marginal cost of overselling; DA15 mechanically
closes the oversell and flips the sign.

The empirical picture in the Spanish data, using OMIE's firm
classification and per-unit repositioning $\Delta Q$, departs from the
model's structural prose in two respects. First, the ISP15 channel
(R1$\to$R2) barely fires in isolation --- the compression at 2024-12-01
is essentially nil once wind is held constant --- which we read as
dominants bidding far from the penalty-binding region, so that a change
in $\kappa$ alone is a second-order shift in their DA programme.
Second, the bulk of the compression falls at R2$\to$R3 (MTU15 intraday,
2025-03-19), where the thinness channel and the penalty channel
activate jointly: within-hour asymmetry reduces $\Phi(\beta)$ and the
firm loses the coarse-granularity hedge that kept the ISP15 penalty
slack under DA60/ID60. The thinness channel is therefore the
load-bearing empirical channel in Spain, not the ISP-window channel in
isolation.

The sign flip predicted at R4 (DA15/ID15, post 2025-10-01) is not
observed in the first three months of data: dominant net intraday
positions remain negative (Big-4 $\Delta Q \approx -11{,}500$
MWh/firm-month). This is consistent with either `flip delayed'
(transitional frictions --- unit-by-unit capacity-matching under DA15
is operationally non-trivial and may take several months of learning)
or `flip absent' (the oversell branch remains active for reasons the
model does not capture, e.g.\ reliability-driven capacity withholding
or inherited DA commitment routines). Three months of DA15/ID15 data
cannot distinguish the two; the test should be revisited with at least
six months of post-2025-10-01 evidence.

Coefficients are directly observable from ESIOS (reserve-activation
prices, system imbalance sign) once $b_{21}$ is pinned down from the
base model; the empirical reassessment above concerns which channel
dominates the Spanish compression, not whether the structural mapping
itself survives.

\begin{thebibliography}{9}

\bibitem[Ito and Reguant(2016)]{ito2016}
Ito, K. and Reguant, M. (2016).
Sequential markets, market power, and arbitrage.
\emph{American Economic Review}, 106(7):1921--1957.

\bibitem[Reneses and Centeno(2014)]{reneses2014}
Reneses, J. and Centeno, E. (2014).
Analysis of the imbalance price scheme in the Spanish electricity market:
A wind power test case.
\emph{Energy Policy}, 67:292--302.

\end{thebibliography}

\appendix

\section{Full model and arbitrage cases}
\label{app:model}

\subsection{Firm's optimisation problem}

The dominant firm is a \emph{residual monopolist}: in each market stage it takes the supply of
competitive fringe suppliers as given. Fringe suppliers bid at marginal cost, so the dominant firm
faces a downward-sloping residual demand.

\paragraph{Day-ahead stage.} The firm commits block output $q_{1t}^{\mathrm{phys}}$ at the
uniform DA clearing price $p_1$, which clears the market against fringe supply. Total
DA residual demand (summing across subperiods) is $D_1(p_1) = A - b_1 p_1$, where $A$
is total forecasted demand and fringe supply has slope $b_1$. Under DA60, the hourly
block forces $q_{11}^{\mathrm{phys}} = q_{12}^{\mathrm{phys}} = q_1/2$; under DA15 each
subperiod is chosen independently.

\paragraph{Intraday stage.} Taking $p_1$ and $q_1^{\mathrm{phys}}$ as given, the firm chooses
$q_{2t}$ in each subperiod (negative $q_{2t}$ means buying back). Fringe suppliers adjust along
their supply curve, giving residual demand
\begin{equation}
D_{2t}(p_1, p_{2t}) \;=\; b_{2t}(p_1 - p_{2t}),
\label{eq:resid2}
\end{equation}
with $b_{2t} \leq b_1$, since production is less flexible closer to
delivery. Inverting~(\ref{eq:resid2}) gives
\begin{equation}
p_{2t} \;=\; p_1 - \frac{q_{2t}}{b_{2t}}.
\label{eq:invdem2}
\end{equation}
Combining intraday revenue with the expected imbalance-settlement cost
from Proposition~\ref{prop:mcstar}, the firm's intraday objective is
\begin{equation}
\max_{q_{21},q_{22}} \;\sum_{t=1}^{2} \bigl[p_{2t}\,q_{2t} - \mathrm{MC}^\star(\omega)\,q_{2t}\bigr].
\label{eq:intraday_obj}
\end{equation}
Substitute~(\ref{eq:invdem2}) into the objective and expand the product:
\begin{align*}
\sum_{t=1}^{2}\!\bigl[p_{2t}\,q_{2t}-\mathrm{MC}^\star q_{2t}\bigr]
  &= \sum_{t=1}^{2}\!\Bigl[
      \bigl(p_1 - q_{2t}/b_{2t}\bigr)\,q_{2t}
      - \mathrm{MC}^\star\,q_{2t}
    \Bigr] \\
  &= \sum_{t=1}^{2}\!\Bigl[
      \bigl(p_1-\mathrm{MC}^\star\bigr)\,q_{2t}
      - q_{2t}^{\,2}/b_{2t}
    \Bigr].
\end{align*}
The subperiods decouple. Differentiating with respect to $q_{2t}$ and
setting the derivative to zero:
\[
\bigl(p_1-\mathrm{MC}^\star\bigr) - \frac{2\,q_{2t}}{b_{2t}} \;=\; 0
\qquad\Longrightarrow\qquad
q_{2t} \;=\; \frac{b_{2t}}{2}\bigl(p_1-\mathrm{MC}^\star(\omega)\bigr).
\]
Summing the two first-order conditions and collecting:
\begin{align*}
q_{21}+q_{22}
  &= \tfrac{b_{21}}{2}\bigl(p_1-\mathrm{MC}^\star\bigr)
     + \tfrac{b_{22}}{2}\bigl(p_1-\mathrm{MC}^\star\bigr) \\
  &= \tfrac{b_{21}+b_{22}}{2}\bigl(p_1-\mathrm{MC}^\star\bigr) \\
  &= b_{21}\,\underbrace{\tfrac{1+\beta}{2}}_{\Phi(\beta)}\,
     \bigl(p_1-\mathrm{MC}^\star\bigr),
\qquad \beta \equiv b_{22}/b_{21},
\end{align*}
which is equation~(\ref{eq:xform}) in the main text. A negative net
position ($q_{21}+q_{22}<0$) means the firm repurchases in intraday
what it oversold in DA---the \citet{ito2016} \emph{withholding}
signature.

\paragraph{Day-ahead stage (backward induction).} Anticipating the
intraday response, the firm maximises total two-stage profit. Evaluating
the intraday equilibrium profit at the optimum: since
$p_{2t}^\star - \mathrm{MC}^\star = (p_1-\mathrm{MC}^\star)/2$ and
$q_{2t}^\star = (b_{2t}/2)(p_1-\mathrm{MC}^\star)$,
\[
\sum_{t}\!\bigl(p_{2t}^\star-\mathrm{MC}^\star\bigr)\,q_{2t}^\star
  \;=\; \sum_t \tfrac{b_{2t}}{4}\bigl(p_1-\mathrm{MC}^\star\bigr)^2
  \;=\; \tfrac{b_{21}+b_{22}}{4}\bigl(p_1-\mathrm{MC}^\star\bigr)^2.
\]
Writing $p_1(q_1) = (A-q_1)/b_1$ (inverted DA residual demand), the
firm's reduced-form programme is
\[
\max_{q_1}\;\Pi(q_1)
\;=\; \bigl(p_1(q_1)-c\bigr)\,q_1
  + \tfrac{b_{21}+b_{22}}{4}\bigl(p_1(q_1)-\mathrm{MC}^\star\bigr)^2.
\]
Using $dp_1/dq_1 = -1/b_1$, the first-order condition is
\[
\underbrace{\bigl(p_1-c\bigr) - q_1/b_1}_{\text{DA monopoly margin}}
  \;-\; \underbrace{\tfrac{b_{21}+b_{22}}{2 b_1}\bigl(p_1-\mathrm{MC}^\star\bigr)}_{\text{intraday-profit response}}
  \;=\; 0.
\]
With linear residual demand and constant marginal cost $c$, the
equilibrium satisfies $\mathrm{MC}^\star > p_1 > c$ on the oversell
branch (the analogue of Result~1 in \citealt{ito2016}): the firm
exercises DA markup over $c$, and the imbalance-settlement premium
$\mathrm{MC}^\star - p_1 = \pi\phi^+\kappa(\omega) > 0$ discourages
further IDA sales so $q_{2t}<0$. Substituting
$\mathrm{MC}^\star = p_1 + \pi\phi^+\kappa(\omega)$ into
equation~(\ref{eq:xform}):
\begin{align*}
p_1 - \mathrm{MC}^\star
  &= p_1 - \bigl[p_1 + \pi\phi^+\kappa(\omega)\bigr]
  \;=\; -\,\pi\,\phi^+\,\kappa(\omega),\\[2pt]
q_{21}+q_{22}
  &= b_{21}\,\Phi(\beta)\cdot\bigl(-\pi\phi^+\kappa(\omega)\bigr)
  \;=\; -\,b_{21}\,\Phi(\beta)\,\pi\,\phi^+\,\kappa(\omega) \;<\; 0,
\end{align*}
so the dominant firm is a systematic net buyer in intraday.

\subsection{Fringe arbitrage and capacity constraints}

Following \citet{ito2016}, fringe suppliers can exploit the DA premium by
overselling quantity $s$ in DA and buying it back in intraday. An
arbitrageur shifting $s$ modifies residual demands to
\[
D_1(p_1,s) \;=\; A - b_1 p_1 - s,
\qquad
D_{2t}(p_1,p_{2t},s) \;=\; b_{2t}(p_1-p_{2t}) + s/2,
\]
with $s/2$ split equally across subperiods. Inverting gives
\[
p_1(q_1,s) = \frac{A - q_1 - s}{b_1},
\qquad
p_{2t}(q_{2t},s) = p_1 - \frac{q_{2t} - s/2}{b_{2t}}.
\]

\paragraph{Case 1: No arbitrage ($s=0$).} The dominant firm exploits the
full DA premium. Evaluating equation~(\ref{eq:xform}) on the oversell
branch:
\[
q_{21}+q_{22}
   \;=\; -\,b_{21}\,\Phi(\beta)\,\pi\,\phi^+\,\kappa(\omega),
\]
which is the maximum-magnitude net position in a given regime.

\paragraph{Case 2: Full competitive arbitrage.} With many fringe wind
farms, $s$ rises until the no-arbitrage condition $p_1 = p_{2t}$ holds
for each $t$. Setting $p_{2t} = p_1$ in the inverse demand gives
\[
0 \;=\; -\frac{q_{2t} - s/2}{b_{2t}}
\qquad\Longrightarrow\qquad
q_{2t}^{\mathrm{arb}} \;=\; s/2,
\]
so the arbitrageur's buy-back in each subperiod exactly matches the
dominant firm's IDA sale. The dominant firm's IDA FOC
$q_{2t} = (b_{2t}/2)(p_1-\mathrm{MC}^\star)$ is unchanged, but when the
no-arbitrage condition binds the effective residual demand it faces in
IDA has slope $(b_1+b_{21}+b_{22})$ across the combined market. At
$p_1=p_2$, the firm is indifferent across stages and chooses total
output $q_1 + q_{21} + q_{22}$ to solve a single monopoly problem
against the pooled slope. Substituting into the firm's DA FOC yields a
lower equilibrium $q_1$ --- the firm \emph{withholds more} in DA ---
and a net intraday position
\[
q_{21}+q_{22}
   \;=\; \frac{b_{21}+b_{22}}{b_1+b_{21}+b_{22}}
         \cdot\bigl(\text{pooled markup adjustment}\bigr),
\]
compressed relative to Case~1 but still negative so long as
$\mathrm{MC}^\star > p_1$. This is the mechanism behind Result~2 of
\citet{ito2016}: full arbitrage converges prices, forces the dominant
firm to withhold more in DA, and shrinks total output relative to the
no-arbitrage benchmark.

\paragraph{Case 3: Limited arbitrage ($s \leq K$).} In practice,
arbitrage is constrained by:
\begin{itemize}
\item \emph{Physical capacity}: fringe firms cannot sell beyond
  installed capacity in DA.
\item \emph{Regulatory limits}: OMIE discourages systematic large
  swings in scheduled production unless justified by technical
  operations \citep{ito2016}.
\item \emph{MTU granularity}: under MTU15, there are four delivery
  subperiods per hour instead of one, multiplying the number of
  $(s,\pi)$ slots the arbitrageur must manage and potentially
  increasing coordination costs.
\end{itemize}
When $s=K$ binds, the arbitrageur's residual-demand shift is capped.
Taking the DA-FOC and IDA-FOC at $s=K$, the DA--IDA wedge solves
\[
p_1 - p_{2t}
   \;=\; \frac{q_{2t} - K/2}{b_{2t}},
\]
so for any $q_{2t}<K/2$ (the dominant firm under-arbitrages relative to
the fringe capacity) the wedge is strictly positive. Following
Result~3 of \citet{ito2016}, substituting the dominant firm's FOC and
solving for the equilibrium residual premium gives comparative statics:
\[
\frac{\partial(p_1-p_2)}{\partial A}>0,\qquad
\frac{\partial(p_1-p_2)}{\partial b_1}<0,\qquad
\frac{\partial(p_1-p_2)}{\partial b_{2t}}<0,
\]
i.e.\ DA premia persist and are larger when forecasted demand $A$ is
high, fringe supply is inelastic (low $b_1$), or the intraday market is
elastic (high $b_{2t}$). All three comparative statics carry directly
into the reform narrative.

\subsection{Regime map and the four channels}

The four Spanish regimes map onto the Ito--Reguant framework as follows.

\begin{description}
\item[R1 (DA60/ID60, ISP60).]
  Within-hour netting ($\kappa_{60}<1$) lowers the imbalance cost of a given DA
  overcommitment, so $\mathrm{MC}^\star$ is close to $p_1$ and the withholding
  incentive is moderated. Fringe wind farms arbitrage the DA premium; the
  arbitrage capacity constraint $K$ is typically binding, leaving a residual
  premium. Net intraday position: $q_{21}+q_{22}<0$.

\item[R2 (DA60/ID60, ISP15).]
  ISP15 sets $\kappa=1$, strictly raising $\mathrm{MC}^\star = p_1+\pi\phi^+$.
  The dominant firm internalises the higher cost at the DA stage and reduces $q_1$,
  compressing $|q_{21}+q_{22}|$ through equation~(\ref{eq:xform}) without a sign
  flip. Fringe arbitrage capacity is unchanged, but it operates against a smaller
  residual premium. Net position: $q_{21}+q_{22}<0$, smaller in magnitude.

\item[R3 (DA60/ID15, ISP15).]
  MTU15 intraday introduces within-hour asymmetry: the second subperiod's residual
  demand slope falls ($b_{22} < b_{21}$, i.e.\ $\beta < 1$), lowering $\Phi(\beta)$
  and further compressing $|q_{21}+q_{22}|$ without changing its sign. The
  dominant firm does not change branches; it simply faces a thinner intraday market
  for repositioning. Net position: $q_{21}+q_{22}<0$, further compressed.

\item[R4 (DA15/ID15, ISP15).]
  DA15 allows the firm to set $q_{1t}^{\mathrm{phys}} = K_t$ slot-by-slot, closing
  the structural oversell. The firm moves to the undersell branch:
  $\mathrm{MC}^\star = p_1-(1-\pi)\phi^- < p_1$, so $p_1-\mathrm{MC}^\star > 0$
  and $q_{21}+q_{22}>0$. The dominant firm becomes a \emph{net seller} in intraday,
  and fringe firms that previously arbitraged by buying back in IDA now face the
  opposite premium structure. Net position: $q_{21}+q_{22}>0$ (sign flip).
\end{description}

The regime table in Section~\ref{sec:regimes} summarises these four cases. The
sign sequence $(-,-,-,+)$ is generated entirely by the two primitives (physical
capacity and dual-priced settlement) interacting with three institutional
changes; it does not require calibration of $\phi^\pm$, $\pi$, or $b_{2t}$.

\end{document}
